\section{Introduction}\label{sec:intro}
Let $-D<0$ be a fundamental discriminant, let $F_D=\mathbb Q(\sqrt{-D})$,
and write $\OO_D$ for its ring of integers. Gauss's class number one problem
asks when $\OO_D$ has unique factorization, or equivalently when its ideal
class number $h(-D)$ is one. Heegner~\cite{heegner}, Baker~\cite{baker} and
Stark~\cite{stark} proved that this occurs exactly for
\begingroup
\abovedisplayskip=14pt
\belowdisplayskip=14pt
\abovedisplayshortskip=14pt
\belowdisplayshortskip=14pt
\[
 D\in\{3,4,7,8,11,19,43,67,163\}.
\]
\endgroup
Here we prove a weaker bound by a different method.

\begin{theorem}\label{thm:main}
Let $-D<0$ be a fundamental discriminant with $h(-D)=1$. Then $D\leq498$.
\end{theorem}

The challenge is to obtain as small an explicit upper bound as possible for
a hypothetical tenth discriminant, without relying on Heegner's methods or
the known classification~\cite{tao2026}. We approach this problem through the
determinant method developed in recent work on
Landau--Siegel zeros~\cite{green2026,taosketch}. In particular, the bound in
Theorem~\ref{thm:main} is obtained from determinant inequalities; we do not
exclude the remaining discriminants by a finite test of quadratic residues.

\subsection{Overview of the proof}\label{subsec:overview}
Suppose that $h(-D)=1$ with $D>498$. We will form a matrix by evaluating
polynomials in $(z,w)$ and their $z$-derivatives at pairs $(z,\bar z)$ with
$z\in\OO_D$. If this
matrix has full row rank, it has a nonzero square minor $\Delta$. Since its
entries lie in $\OO_D$, the norm $N=\Delta\bar\Delta=|\Delta|^2$ is a positive
integer. We will bound $N$ from above using the absolute values of the
entries, and show that it is divisible by an integer larger than this upper
bound. This contradiction will exclude the proposed value of $D$. The
choice of polynomials and evaluation points must therefore serve both
estimates while ensuring that the matrix has full row rank.

To see how the class number assumption enters, suppose that a prime $p$ is
not inert. There is then an ideal of norm $p$, which must be principal because
$h(-D)=1$. Its generator has norm $p$ and cannot be a rational integer. But
every non-rational element of $\OO_D$ has norm at least $D/4$, so every prime
$p<D/4$ must be inert. At such a prime, conjugation agrees with Frobenius
modulo $p$, giving $z^p\equiv\bar z\pmod{p\OO_D}$. Thus the polynomial
$z^p-w$ takes values divisible by $p$ at every point $(z,\bar z)$ with
$z\in\OO_D$. By expressing the rows of our matrix in terms of powers of
$z^p-w$, we obtain prime powers dividing $N$.

This change of basis must remain within the chosen polynomial space. We use
the monomials $z^rw^s$ with $r+Ks<R$, where $K$ and $R$ are parameters. For
$p\geq K$, replacing $z^p$ by $w$ does not increase $r+Ks$, so the triangular
space of exponents permits all the required row operations. Taking several derivatives at each point gives more columns from the same
set of lattice points. Differentiation reduces the resulting divisibility,
but by an amount that can be bounded explicitly in terms of the derivative
order. We also use inert primes above
$D/4$: the same norm argument bounds the number of non-inert primes in an
interval by the number of elements of $\OO_D$ whose norms lie there.

For the upper bound, we choose the $n$ lattice points of $\OO_D$ nearest to
the origin. Reducing $n$ reduces the radius of the disk containing them and
hence the bounds on the polynomial values. The estimates must also account
for the derivative columns at each point. These columns are far from
orthogonal, so bounding each separately loses too much at the smaller
values of $D$. Instead, we estimate their Gram matrix, retaining the saving
from these correlations before combining the contributions of the
different points.

It remains to determine how few points we can use without losing full row
rank. If the polynomial space has dimension $M$ and we take $J$ derivative
orders at each point, there are $nJ$ columns, so necessarily $nJ\geq M$.
We prove a sufficient condition close to this necessary one. Direct
reduction modulo an auxiliary inert prime $p$ loses information because
$z^p-w$ vanishes at every reduced node. We therefore rescale both the
polynomial rows and the derivative columns using their known divisibility,
obtaining an integral matrix whose reduction retains this information.
A hypothetical row relation then gives a polynomial in two variables with
prescribed local vanishing. Applying a differential operator and taking a
determinant of coefficients produces a nonzero polynomial in one variable
whose zeros, counted with multiplicity, exceed its degree. This contradiction
permits a smaller $n$ and therefore improves the upper bound for $N$.

We apply these estimates as follows. Elementary norm arguments first reduce
$D>498$ to primes congruent to $3$ modulo $8$. A uniform estimate and $101$
adjacent interval comparisons exclude all $D\geq875$, leaving sixteen
possible values below $875$. Fifteen are excluded by fixed determinant
comparisons; the last is $499$. Table~\ref{tab:opening} shows the effect of
the smaller node count there, with the polynomial space and derivative
orders fixed. After the derivative-order contributions have been combined,
the arithmetic term changes only slightly, whereas the size estimate
decreases enough to make the normalized difference positive.

\begin{table}[ht]
\centering
\begin{tabular}{lrr}
\toprule
 & Larger node count & Smaller node count\\
\midrule
Nodes $n$ & $41\,793\,600$ & $35\,394\,400$\\
Proportion $M/(nJ)$ & $0.83333\ldots$ & $0.98399\ldots$\\
Arithmetic minus size bound & $-0.12343\ldots$ & $+0.05119\ldots$\\
\bottomrule
\end{tabular}
\caption{The comparison at $D=499$, with the polynomial space and derivative
orders fixed (approximate values). The directed bounds are given in
Section~\ref{sec:completion}.}
\label{tab:opening}
\end{table}

Section~\ref{sec:comparison} introduces the determinant and the inequality
we must contradict. Sections~\ref{sec:arithmetic} and~\ref{sec:metric} develop
its divisibility and size estimates. Section~\ref{sec:rank} proves the rank
criterion, and Section~\ref{sec:completion} completes the numerical argument.
The explicit arithmetic formulas, subsidiary spectral proofs, uniform
estimates for $D\geq875$, and numerical calculations are collected in the
accompanying repository,
\url{https://github.com/dualverse-ai/tenth_discriminant}.

\section{The determinant comparison}\label{sec:comparison}

\subsection{A model: an integer that is small and highly divisible}
\label{subsec:model}

Every nonzero integer $N$ satisfies
\begin{equation}\label{eq:product-formula}
  \log|N|=\sum_p v_p(N)\log p .
\end{equation}
Our integer will be $N=|\Delta|^2$ for a nonzero algebraic integer $\Delta\in
\mathcal{O}_D$, which is the norm of $\Delta$ and hence a positive rational
integer. The strategy is to bound $\log N$ from above by Archimedean
estimates, to bound $v_p(N)$ from below at many primes, and to show that the
two bounds are incompatible when $h(-D)=1$. This subsection illustrates the comparison using elementary estimates. The rest of the paper sharpens each of its
three ingredients.

\subsubsection*{The matrix}
We work with polynomials in two independent variables $z$ and $w$. The variable
$w$ will be specialized to the complex conjugate of the value of $z$, but it is
treated as independent of $z$ throughout. For $P\in\mathbb{Z}[z,w]$ define the
Hasse derivatives in $z$, with $w$ held fixed, by
\[
  P(z+t,w)=\sum_{j\ge0}\bigl(\partial_z^{[j]}P\bigr)(z,w)\,t^j,
  \qquad
  \partial_z^{[j]}\bigl(z^rw^s\bigr)=\binom rj z^{r-j}w^s .
\]
Hasse derivatives keep integer coefficients; ordinary derivatives would
introduce a factor $j!$. Fix a positive integer $J$. Let $z_1,\dots,z_n\in\mathcal{O}_D$ be
distinct \emph{nodes}, and let $\Lambda\subset\mathbb{Z}_{\ge0}^2$ be a finite
set of exponents. The \emph{jet matrix} has rows indexed by $(r,s)\in\Lambda$,
columns indexed by pairs $(i,j)$ with $1\le i\le n$ and $0\le j<J$, and entries
\begin{equation}\label{eq:jet-matrix}
  V_{(r,s),(i,j)}
  =\bigl(\partial_z^{[j]}z^rw^s\bigr)(z_i,\overline{z_i})
  =\binom rjz_i^{\,r-j}\,\overline{z_i}^{\,s}\in\mathcal{O}_D .
\end{equation}
For a set $C$ of $|\Lambda|$ columns, write $\Delta_C$ for the corresponding
maximal minor. More generally, the row of any polynomial $P$ in the span of the
monomials $z^rw^s$, $(r,s)\in\Lambda$, is the vector of its jets
$(\partial_z^{[j]}P)(z_i,\overline{z_i})$. Integral unitriangular
changes of the spanning polynomials leave every $\Delta_C$ unchanged.

\subsubsection*{One Frobenius calculation}
The divisibility comes from the behaviour of conjugation modulo an inert prime.

\begin{lemma}\label{lem:frobenius}
Let $p$ be a prime that is inert in $F_D$, and let $z\in
\mathcal{O}_D$. Then $z^p\equiv\overline z\pmod{p\mathcal{O}_D}$.
Consequently, if $g=z^p-w$, then
\[
  g(z+t,\overline z)=pU(t)+t^p
\]
for a polynomial $U\in\mathcal{O}_D[t]$ of degree less than $p$, and for all
integers $m,j\ge0$,
\begin{equation}\label{eq:layer-divisibility}
  p^{\lfloor j/p\rfloor}\,[t^j]\,g(z+t,\overline z)^m\in
  p^m\mathcal{O}_D .
\end{equation}
\end{lemma}

\begin{proof}
The quotient $\OO_D/p\OO_D$ is the field of $p^2$ elements. Conjugation
induces its nontrivial automorphism, which is the $p$th-power map. This gives
$z^p\equiv\bar z\pmod{p\OO_D}$, including when $p=2$. Next,
\[
  (z+t)^p-\overline z
  =(z^p-\overline z)+\sum_{k=1}^{p-1}\binom pkz^{p-k}t^k+t^p,
\]
and every term except $t^p$ is divisible by $p$. Finally,
$(pU+t^p)^m=\sum_{q=0}^m\binom mq p^{m-q}U^{m-q}t^{pq}$, and a term contributes
to $t^j$ only if $q\le\lfloor j/p\rfloor$. Each such term is divisible by
$p^{m-\lfloor j/p\rfloor}$, which proves~\eqref{eq:layer-divisibility}.
\end{proof}

Geometrically, the lemma says that, $p$-adically, every point
$(z,\overline z)$ lies close to the curve $z^p=w$~\cite{taosketch}.
The polynomial $g$ is small at every node, and its $z$-derivatives are small
too, since $\partial_z g=pz^{p-1}$. This is why the jets differentiate only in
$z$: a derivative in $w$ would not preserve the divisibility.

\subsubsection*{Where the triangle comes from}
To use Lemma~\ref{lem:frobenius}, write each row exponent as $r=mp+a$ with
$0\le a<p$ and replace the monomial $z^rw^s$ by
\begin{equation}\label{eq:frobenius-basis}
  g^m z^aw^s=z^rw^s+\sum_{i=1}^m(-1)^i\binom mi z^{r-pi}w^{s+i}.
\end{equation}
By the Leibniz rule and~\eqref{eq:layer-divisibility}, every jet of
$g^mz^aw^s$ of order $j$, multiplied by $p^{\lfloor j/p\rfloor}$, lies in
$p^m\mathcal{O}_D$. The change of basis~\eqref{eq:frobenius-basis} is
unitriangular for the ordering by $z$-degree, but only if every exponent
$(r-pi,s+i)$ on the right is again in $\Lambda$. The support must therefore be
closed under the Frobenius substitution $z^p\mapsto w$.

The rectangle $\{r<R,\ s<S\}$ used by Green is not closed under this
substitution: it is stopped by the ceiling $s<S$. There the substitution can be
applied only $\min(\lfloor r/p\rfloor,S-1-s)$ times, and the $p$-adic credit of
the row is truncated accordingly~\cite[Lemma~5.1]{green2026}. The remedy is to
grade the monomials so that the substitution never raises the degree. Fix integers $K\geq2$ and $S\geq1$. Give $z$
weight $1$ and $w$ weight $K$, and take
\begin{equation}\label{eq:triangle}
  \Lambda=\{(r,s)\in\mathbb{Z}_{\ge0}^2:\ r+Ks<R\},\qquad R=KS,\qquad
  M=|\Lambda|=\frac{R(S+1)}{2}.
\end{equation}
The substitution $z^p\mapsto w$ changes the weight by $K-p$, so it never leaves
$\Lambda$ when $p\ge K$. Every row then receives its full credit
$\lfloor r/p\rfloor$. Primes $p<K$ receive truncated credit, and that boundary
is treated in Section~\ref{sec:arithmetic}.

Scaling column $(i,j)$ by $p^{\lfloor j/p\rfloor}$ multiplies $\Delta_C$ by
$p^{\sum_{(i,j)\in C}\lfloor j/p\rfloor}$. Since $p\mathcal{O}_D$ is prime and
$v_p(|\Delta_C|^2)=2v_{p\mathcal{O}_D}(\Delta_C)$, we obtain, for every inert
prime $p\geq K$,
\begin{equation}\label{eq:model-valuation}
  v_p\bigl(|\Delta_C|^2\bigr)\ \ge\
  2\sum_{(r,s)\in\Lambda}\Bigl\lfloor\frac rp\Bigr\rfloor
  \;-\;2\sum_{(i,j)\in C}\Bigl\lfloor\frac jp\Bigr\rfloor .
\end{equation}
The row contribution, including its factor $2$, is about $W/p$, where
$W=R^2S/3\approx2\sum_\Lambda r$. The second is a column loss that is present
only when derivative orders reach $p$. In our application $J$ is much larger
than $p$, so the loss is real. Section~\ref{sec:arithmetic} shows how it
combines with the factorial sizes of the Hasse coefficients.

\subsubsection*{Sizes}
If every node satisfies $|z_i|\le\rho$ with $\rho\ge1$, then each entry
of row $(r,s)$ has absolute value at most $\binom rj\rho^{r+s-j}$. Hadamard's
inequality therefore gives
\begin{equation}\label{eq:model-hadamard}
  \log|\Delta_C|^2\ \le\ \sum_{(r,s)\in\Lambda}2(r+s)\log\rho
  \;+\;\sum_{(r,s)\in\Lambda}\log\sum_{(i,j)\in C}\binom rj^2\rho^{-2j}.
\end{equation}
The first term is $(1+1/K+O(1/S))\,W\log\rho$. Combining
\eqref{eq:product-formula}, \eqref{eq:model-valuation} and
\eqref{eq:model-hadamard}, any nonzero $\Delta_C$ satisfies
\begin{equation}\label{eq:model-inequality}
  \sum_{\substack{p\geq K\\ p\text{ inert}}}2\Bigl(\sum_{\Lambda}
  \Bigl\lfloor\frac rp\Bigr\rfloor-\sum_{C}\Bigl\lfloor\frac jp\Bigr\rfloor
  \Bigr)\log p
  \le\ \sum_{\Lambda}2(r+s)\log\rho
  +\sum_{\Lambda}\log\sum_{C}\binom rj^2\rho^{-2j}.
\end{equation}
After division by $W$, the left side is roughly the sum of $(\log p)/p$ over
inert primes $K\leq p<R$, less the column loss. Under $h(-D)=1$ this inert sum is
bounded below by the corresponding sum over all primes, minus a loss controlled
by counting values of the norm form, as in Section~\ref{sec:arithmetic}. The
right side is roughly $(1+1/K)\log\rho$ plus the derivative term. The theorem
follows when the left side is larger.

This model already shows where the difficulty lies, and it points to the three
refinements. First, the arithmetic credit must be computed exactly, including
small primes, prime powers, and the loss from non-inert primes, bounded through
the norm form (Section~\ref{sec:arithmetic}). Second, Hadamard's inequality treats the
entries of each node block as unrelated and discards their cancellation. The
spectral estimates of Section~\ref{sec:metric} recover part of it. Because the
nonzero minor is not identified, those estimates must hold for every column set
$C$. Third, $\rho$ is governed by the node count. The $n$ nearest lattice points
lie in a disk of radius about $\sqrt{n\sqrt D/(2\pi)}$
(Lemma~\ref{lem:node-disk}), and $n\ge M/J$ is forced.

\subsection{Arithmetic of class number one}\label{subsec:arith}

Every element of $\mathcal{O}_D$ has the form $\alpha=(x+y\sqrt{-D})/2$ with
$x,y\in\mathbb{Z}$ and $x\equiv yD\pmod 2$, and its norm is
\begin{equation}\label{eq:norm}
  N(\alpha)=|\alpha|^2=\frac{x^2+Dy^2}{4}.
\end{equation}
The elementary reductions below follow from one observation about norms.

\begin{lemma}[Norm obstruction]\label{lem:norm-obstruction}
If $\alpha\in\mathcal{O}_D\setminus\mathbb{Z}$, then $N(\alpha)\ge D/4$. If
$\alpha\in\mathbb{Z}$, then $N(\alpha)$ is a square.
\end{lemma}

\begin{proof}
If $\alpha\notin\mathbb{Z}$ then $y\ne0$ in~\eqref{eq:norm}, so $N(\alpha)\ge
Dy^2/4\ge D/4$. If $\alpha\in\mathbb{Z}$ then $N(\alpha)=\alpha^2$.
\end{proof}

\begin{lemma}\label{lem:small-inert}
Suppose $h(-D)=1$. Then every prime $p<D/4$ is inert in $F_D$.
\end{lemma}

\begin{proof}
If $p$ is split or ramified, there is a prime ideal of norm $p$. It is
principal, generated by some $\alpha$ with $N(\alpha)=p$. Since $p$ is not a
square, $\alpha\notin\mathbb{Z}$, and Lemma~\ref{lem:norm-obstruction} gives
$p\ge D/4$.
\end{proof}

\begin{lemma}[Reduction]\label{lem:reduction}
Suppose $h(-D)=1$ and $D>16$. Then $D$ is a prime and $D\equiv3\pmod 8$.
\end{lemma}

\begin{proof}
The ramified primes are the prime divisors of $D$, so by
Lemma~\ref{lem:small-inert} every prime divisor of $D$ is at least $D/4$. If
$D$ is even, then $2$ ramifies and $2<D/4$, which is impossible. Hence $D$ is
odd, and since $-D$ is a fundamental discriminant, $D$ is squarefree with
$D\equiv3\pmod4$. If $D$ were composite, its least prime factor $q$ would
satisfy $q\le\sqrt D<D/4$, again impossible. So $D$ is prime. Finally, if
$D\equiv7\pmod8$ then $-D\equiv1\pmod 8$ and $2$ splits, contradicting
Lemma~\ref{lem:small-inert}. Therefore $D\equiv3\pmod8$.
\end{proof}

For the proof of Theorem~\ref{thm:main} we assume
$D>498$ from now on, so $D$ is a prime congruent to $3$ modulo $8$, and both $2$
and every odd prime below $D/4$ are inert. Since $D\equiv3\pmod4$, the
integer cutoff $p<D/4$ is equivalently $p\leq(D-3)/4$. We write
$\omega=(1+\sqrt{-D})/2$, so that $\OO_D=\ZZ[\omega]$. In these coordinates
$\N(x_0+y\omega)=(x_0+y/2)^2+Dy^2/4$; it is the same norm form as
\eqref{eq:norm}, with $x=2x_0+y$.

\subsection{Nearest nodes}\label{subsec:geometry}

Our nodes are the lattice points of $\mathcal{O}_D\subset\mathbb{C}$ nearest to
the origin. Let $z_1,z_2,\dots$ be the elements of $\mathcal{O}_D$
listed with $|z_1|\le|z_2|\le\cdots$, with ties broken arbitrarily.
By~\eqref{eq:norm}, $\mathcal{O}_D$ is the union of the horizontal rows
\[
  \Bigl\{\tfrac{x}{2}+\tfrac{y}{2}\sqrt{-D}:\ x\equiv yD\ (\mathrm{mod}\ 2)\Bigr\},
  \qquad y\in\mathbb{Z}.
\]
The rows are spaced $\sqrt D/2$ apart vertically, and within each row the
points form a translate of $\mathbb{Z}$. The covolume of $\mathcal{O}_D$ is
therefore $\sqrt D/2$.

\begin{lemma}[Nearest-node disk]\label{lem:node-disk}
Put $\delta=\sqrt{D+4}/4$. For every $i\ge1$,
\begin{equation}\label{eq:node-disk}
  |z_i|\ \le\ \rho_i:=\sqrt{\frac{i\sqrt D}{2\pi}}+\delta .
\end{equation}
\end{lemma}

\begin{proof}
Every point of the plane is within vertical distance $\sqrt D/4$ of a
lattice row, and within horizontal distance $1/2$ of a point on that row.
Thus the lattice has covering radius at most $\delta$.
For $z\in\mathcal{O}_D$, let $V(z)$ be its Voronoi cell, the set of
points of $\mathbb{C}$ at least as close to $z$ as to any other lattice
point. These cells are translates of one another, they cover $\mathbb{C}$, and
they overlap only on their boundaries. Each has area $\sqrt D/2$. Every point of $V(z)$ lies within $\delta$ of $z$.

Let $\rho>\delta$. If $|\zeta|\le\rho-\delta$ and $\zeta\in V(z)$, then
$|z|\le|\zeta|+\delta\le\rho$. Hence the closed disk of radius
$\rho-\delta$ is covered by the cells of lattice points of modulus at most
$\rho$, and comparing areas gives
\[
  \#\{z\in\mathcal{O}_D:\ |z|\le\rho\}\ \ge\
  \frac{2\pi(\rho-\delta)^2}{\sqrt D}.
\]
With $\rho=\rho_i$ the right side equals $i$, so at least $i$ lattice points
have modulus at most $\rho_i$. This gives $|z_i|\le\rho_i$.
\end{proof}

We use Lemma~\ref{lem:node-disk} in two ways. It supplies the radius
$\rho=\rho_n$ in the Archimedean estimates. It also controls residues, through
the following consequence.

\begin{corollary}[Distinct residues]\label{cor:distinct-residues}
Let $p$ be a prime with $2\rho_n<p$. Then $z_1,\dots,z_n$ have
distinct residues modulo $p\mathcal{O}_D$.
\end{corollary}

\begin{proof}
A nonzero element $p\beta$ of $p\mathcal{O}_D$ has modulus
$p\sqrt{N(\beta)}\ge p$. The difference of two distinct nodes is nonzero and
has modulus at most $2\rho_n<p$.
\end{proof}

At $D=499$, with $n=35\,394\,400$, Lemma~\ref{lem:node-disk} gives
$2\rho_n<22\,491.5$. This is far below every candidate auxiliary prime in
$(160\,000,170\,000]$, so the residue hypothesis of the rank theorem holds
whichever of those primes is inert.


\subsection{The inequality to be proved}
\label{subsec:central}
We now give the comparison in a form that separates the algebra from the numerical estimates. The notation introduced here will be used throughout the paper. Write
\begin{equation}\label{eq:normalization}
 W=\frac{R^2S}{3},\qquad
 W_1=2\sum_{(r,s)\in\Lambda}r,\qquad
 W_2=2\sum_{(r,s)\in\Lambda}s.
\end{equation}
The factor $W$ is a convenient normalization: the main logarithmic terms in both estimates have size $W\log D$. Direct summation gives
\begin{equation}\label{eq:radius-coefficient}
 \frac{W_1+W_2}{W}
 =1+\frac1K+\frac{3(K-1)}{2KS}+\frac{K-5}{2KS^2}.
\end{equation}

Choose $M$ columns $C$ for which the square submatrix $V_C$ is nonsingular, and put $\Delta_C=\det V_C$. The rank argument will supply such a choice; it will not specify the columns. The arithmetic estimates therefore have to hold for every possible choice $C$. They take the form
\begin{equation}\label{eq:arithmetic-interface}
 W\arith-2\sum_{(i,j)\in C} b(j)\ \leq\ \log|\Delta_C|^2.
\end{equation}
Here $\arith$ is a column-independent lower bound assembled from local divisibilities, and $b(j)$ is the loss from using a derivative of order $j$. Both depend on the chosen parameters. Section~\ref{sec:arithmetic} describes their construction and states the local estimate used below.

The main part of $b(j)$ is the first-prime-power sum
$\sum_q\lfloor j/q\rfloor\log q$. Legendre's formula shows that subtracting
it from $\log(j!)$ leaves higher-prime-power savings, with additional
contributions from non-inert primes. This remaining function is
nearly linear in $j$; we subtract a slope and bound its finite remainder. Set
\begingroup\postdisplaypenalty=10000
\begin{equation}\label{eq:F-H}
 F(j)=\log(j!)-b(j),\qquad H(j)=j\sigma-F(j),
\end{equation}
\endgroup
where $\sigma$ is a slope chosen in the arithmetic calculation. Let $c$ be any bound such that
\begin{equation}\label{eq:column-remainder}
 2\sum_{(i,j)\in C}H(j)\leq Mc
 \quad\hbox{for every }|C|=M.
\end{equation}
There are only $n$ columns of any one order $j$, so this is a finite selection problem with a capacity constraint. It need not be bounded by the much cruder $2M\max_jH(j)$.

Let $\rho$ contain all the nodes, put $u_i=z_i/\rho$, and define
\begin{equation}\label{eq:residual-matrix}
 B_{(r,s),(i,j)}(\ell)
 =\frac{\fall rj}{R^j}e^{\ell j}u_i^{r-j}\bar u_i^{s},
 \qquad \ell=\log(R/\rho)-\sigma.
\end{equation}
Here $\fall rj=r(r-1)\cdots(r-j+1)$ is the falling factorial, with
$\fall r0=1$; the entry is zero when $j>r$. The exact relation to the integral matrix is
\begin{equation}\label{eq:diagonal-scaling}
 B_C(\ell)
 =\operatorname{diag}_{\Lambda}(\rho^{-r-s})\,
 V_C\,
 \operatorname{diag}_{C}\bigl(j!(\rho/R)^je^{\ell j}\bigr).
\end{equation}
Taking squared absolute determinants, substituting \eqref{eq:F-H}, and then using \eqref{eq:arithmetic-interface} gives
\[
 W\arith\leq (W_1+W_2)\log\rho
       +\log|\det B_C(\ell)|^2+2\sum_C H(j).
\]
This calculation explains the normalization: the factorials are cancelled exactly, rather than estimated twice.

\begin{proposition}[Determinant comparison]\label{prop:comparison}
Suppose that $V$ has full row rank, that \eqref{eq:arithmetic-interface} and \eqref{eq:column-remainder} hold, and that
\[
 \frac1W\log|\det B_C(\ell)|^2\leq U_{\rm var}
 \qquad (|C|=M)
\]
holds uniformly over all column selections. Then
\begin{equation}\label{eq:master}
 \arith\leq \frac{W_1+W_2}{W}\log\rho
                  +U_{\rm var}+\frac{Mc}{W}.
\end{equation}
In particular, a strictly larger arithmetic lower bound is impossible.
\end{proposition}
\begin{proof}
Choose a nonsingular minor and apply the preceding identities. All bounds are uniform in the choice, so no information about which minor is nonsingular is needed.
\end{proof}

For the numerical comparisons we use the slightly weaker bound obtained
by adding $2M\log M/W$ to the right side of~\eqref{eq:master}.
This positive term is less than $3\times10^{-7}$ at the parameters for $499$.
Thus define
\begin{equation}\label{eq:metric-total}
 \metric=\frac{W_1+W_2}{W}\log\rho+U_{\rm var}
             +\frac{Mc+2M\log M}{W},
 \qquad G=\arith-\metric.
\end{equation}
A positive lower bound for $G$ rules out the proposed field. There are three separate tasks: prove full rank, establish the local arithmetic estimates, and bound the same minor at the complex embedding. The auxiliary-prime scalings used to prove rank in Section~\ref{sec:rank} do not enter \eqref{eq:diagonal-scaling}; after nonvanishing has been established, all size and divisibility estimates concern the original integral minor $\Delta_C$.

It is already possible to see why the number of nodes matters. The disk estimate gives $\rho\asymp n^{1/2}D^{1/4}$ once its boundary error is small. Reducing $n$ therefore lowers the common term in \eqref{eq:metric-total} by approximately
\[
 \frac12\left(1+\frac1K\right)\log\frac{n_{\rm old}}{n_{\rm new}}.
\]
At the two node counts in Table~\ref{tab:499}, this is about $0.085$. The remaining terms also change: there are fewer copies of each derivative order, and the proportion $\theta=M/(nJ)$ of available columns that must be selected increases. We call this proportion the \emph{occupancy}. Section~\ref{sec:metric} exploits these restrictions. In the actual calculation the combined gain is enough to change the sign of $G$.

